\chapter{„עשרים שניות”}
\label{sec:1}
קצין מודיעין, שתיאר את תפקידו בפס ייצור של מטרות שנבנה סביב רשימות חיסול שהפיקה מכונה, סיפר לעיתונאי שהקדיש עשרים שניות לכל שם, מספיק כדי לוודא שהמטרה היא גבר: „לא לוקח הרבה זמן להבחין אם למישהו יש קול של גבר או של אישה”. הוא הוסיף שלא היה לו שום ערך מוסף כבן אדם, מלבד היותו חותמת אישור. „זה חסך הרבה זמן” \autocite{abraham2024lavender}. לפי עדותו, לזה הסתכם „אדם בלולאה” תחת קצב מבצעי.\footnote{מחקר על הטיית אוטומציה צופה את ההאצלה הזאת: כשדרישות אחרות מתחרות על תשומת הלב שלהם, אנשים שמקבלים המלצה של מכונה בודקים אותה פחות ונשענים עליה יותר \autocite{parasuraman2010complacency}.}

פרשנות אחת היא שהקצין הוא סובייקט שקיבל מעט מדי זמן, והתיקון הוא למשול בו טוב יותר – יותר זמן, בוחן שני, עוד מצפון שמותקן בלולאה. דרך העדשה הזאת של \emph{הכפפה חברתית}, האדם הוא סובייקט שעומד מחוץ למכונה ומפעיל אותה ככלי: נהג ומכונית, פועל ומחרטה. לאדם יש תפקיד, שם, אחריות ומצפון שאפשר לפנות אליו. עדשה חלופית היא \emph{שעבוד מכונתי}, שבו האדם הוא רכיב של המכונה, מחובר עם חלקיה האחרים למערכת שמווסתת את כולם יחד; המכונה משתמשת באדם כאחד מאיבריה, ומה שהאדם תורם הוא פונקציה שמבוצעת בקצב שהמערכת קובעת.\footnote{השם שדוקטורוב נותן לעמדה של הקצין הוא הקנטאור ההפוך: לא אדם שמכונה מסייעת לו, אלא רכיב של מכונה שעשוי מאדם \autocite{doctorow2026reversecentaur}.} השתיים אינן סותרות זו את זו; אדם יכול למלא תפקיד ברובד אחד ובה בעת להיות חלק ברובד האחר \autocite[מישור 13, „The Apparatus of Capture”: יש שעבוד כאשר בני אדם הם עצמם חלקים מכוננים של מכונה שהם מרכיבים בינם לבין עצמם ועם דברים אחרים, תחת שליטתה והכוונתה של אחדות עליונה, ויש הכפפה כאשר האחדות העליונה מכוננת את האדם כסובייקט הקשור לאובייקט חיצוני חדש, שיכול להיות בעל חיים, כלי או אפילו מכונה]{deleuze1987plateaus}, והקצין של עשרים השניות היה שניהם. כחלק, הוא היה המקום שבו פס הייצור הפיק את החותמת שלו, ועשרים שניות היו זמן המחזור של החלק הזה. כסובייקט, פנו אליו לאורך כל הדרך: הוא תיאר לעיתונאי את הדיסוננס בין הידיעות על מות אזרחים לבין הלחץ בעבודה להביא עוד מטרות. וגם כשהדיסוננס הגיע למצפונו, הוא הלך עם הזרם. זה נקלט ולא הכריע דבר;\footnote{או, לפחות, שום דבר שמנע מצבא ההגנה לישראל להרוג 100,000 בני אדם בעזה בין אוקטובר 2023 לאוקטובר 2025 \autocite{gomezugarte2025uncertainty,acosta2026gaza,economist2025gaza}. והעדות של הקצין המסוים הזה פורסמה באפריל 2024 ב־+972 Magazine. אפשר רק לנחש אם ההחלטה של הקצין לדבר עם עיתונאי מנעה מוות נוסף אחד, 50,000 מקרי מוות נוספים, או אף לא אחד. ומה היה קורה אילו הקצין היה מסרב אחרי שבועיים בתפקיד – נניח, באופן היפותטי, בנובמבר 2023? מדובר בקצין מודיעין שהשתמש ב־Lavender כדי לערוך רשימות חיסול. סביבת העבודה היומיומית שלו הייתה כנראה קרובה יותר לעבודה משרדית אזרחית מאשר לסרט פעולה. האם היה עליו לבקש פגישות חוזרות ונשנות עם הממונים עליו, במעלה שרשרת הפיקוד, ובינתיים לסרב לאשר מטרות לתקיפות אוויריות שהיה סביר לצפות שיגרמו למות אזרחים? אפשר לטעון שזה היה הדבר האתי לעשות. בתרחיש הנגדי הזה, האם היו פחות מתים?} את המסקנה שהמדיניות הייתה בלתי מידתית, כך אמר, הוא הסיק בדיעבד.

אז מה שחסר לא היה מצפון; היה לו מצפון. מצפון פועל בכך שהוא גורם להחלטה לעלות משהו, אבל כאן המחירים התקזזו בכיוון ההפוך: „לחצו עלינו כל הזמן: ‚תביאו לנו עוד מטרות.’ ממש צעקו עלינו” \autocite{abraham2024lavender}. הוא היה רכיב עם נקיפות מצפון, אבל הקצב נשמר: נקיפות המצפון לא שינו דבר במה שהוא תרם, והוא הפך לחותמת. רכיב עם נקיפות מצפון שלא עומד בקצב הופך ל„צוואר הבקבוק האנושי” שמפקד לשעבר של יחידת מודיעין האותות של ישראל כתב ספר על הסרתו \autocite{ys2021humanmachine}. מה שהפך את הקצין לרכיב היה ההסדר, לא הגישה שלו. רכיב יכול ללכת עם הזרם בחוסר רצון או בלהיטות: הפקיד חסר המחשבה שארנדט ראתה במשפט אייכמן \autocite{arendt1963eichmann} היה, על פי מה שאמר בארגנטינה ארבע שנים קודם לכן, כשלא היה מונח על הכף דבר, הצגה שהעלה בשביל ירושלים \autocite{stangneth2014eichmann}. מה שפס הייצור איפשר לקצין לתרום היה פעולה קבועה בקצב קבוע: הוא ראה כל מקרה כפי שפס הייצור תיאר אותו, אחד בכל פעם, בזמן שפס הייצור נתן לו. השאלה אם משהו יכול לעמוד בעמדה הזאת ולא להיות רכיב היא השאלה של החלק הרביעי.

תחקירים עיתונאיים תיארו שלוש מערכות שפעלו יחד בעזה: אחת שהפיקה מטרות של מבנים ותשתיות, כשמקורות אמרו שרוב ההצעות היו בתים פרטיים; שנייה שדירגה אנשים לפי השתייכות משוערת לארגונים חמושים, והפיקה רשימות שלפי הדיווח סימנו עד 37,000 בני אדם, ושבפנים הובן שהן שגויות בערך פעם אחת מכל עשר\footnote{הצבא הישראלי חלק על הדיווח. גם הדיווחים על מינאב ועל Maven שלהלן נשענים על תחקירים עיתונאיים שהצדדים חולקים עליהם. אני מביא אותם כפי שדווחו.}; ושלישית שעקבה מתי אנשים שסומנו חזרו הביתה, כדי שאפשר יהיה לתזמן תקיפות אוויריות לשעה שבה היו עם משפחותיהם.

\runin{Maven: הגרסה האמריקאית, עכשיו עם האצת GenAI}

ההפצצה של בית הספר במינאב הרגה 156 בני אדם, 120 מהם תלמידי בית ספר \autocite[פס׳~92]{UNFFMIran2026}.\footnote{הספירות משתנות לפי מקור ותאריך, מ־156 (הספירה הסופית של התובע במינאב, שגם אמנסטי אינטרנשיונל מביא) ועד 168 (ראש העיר, בימים הראשונים). משלחת האו״ם מביאה את נתוני התובע, ומדווחת גם על מידע עצמאי על 157 הרוגים, 123 מהם ילדים בני 13 ומטה \autocite[פס׳~92–93]{UNFFMIran2026}.} התפקיד האזרחי של בית הספר נראה בצילומי לוויין ממקורות גלויים, וקידוד המטרה עדיין נשא ייעוד שעבר בירושה ממתקן סמוך של משמרות המהפכה. משלחת של האו״ם לבירור עובדות דיווחה שיש יסוד סביר להאמין שההפצצה הייתה פשע המלחמה של הפיכת אובייקט אזרחי למושא של מתקפה \autocite{ohchr2026minabpress}, ושבית הספר היה ניתן לזיהוי ברור, בלי שום מידע שמעיד על שימוש צבאי \autocite[פס׳~81–86, 124–126]{UNFFMIran2026}. היא קבעה שארצות הברית לא עמדה בחובתה לעשות כל הניתן באופן מעשי כדי לאמת את המטרה, מנתה את עדכון נתוני המטרות בין אמצעי הזהירות הניתנים לביצוע, וקבעה שהכשל חרג מרשלנות \autocite[פס׳~124–125]{UNFFMIran2026}. אמנסטי אינטרנשיונל הגיע למסקנה מקבילה, בלשון שתואמת את חובת האימות שבסעיף 57(2)(a)(i) לפרוטוקול הנוסף הראשון: שייתכן שהכוחות הסתמכו על מודיעין מיושן ולא עשו כל הניתן באופן מעשי כדי לאמת את המטרה \autocite{Amnesty2026Minab}.

זה קרה בבוקר 28 בפברואר, היום הראשון של המערכה באיראן \autocite{Amnesty2026Minab,crs2026iranwar}. יום קודם לכן הודיע השר שהספק (Anthropic) שהמודל שלו (Claude) ישב בתוך פלטפורמת המטרות (Maven Smart System) יוגדר סיכון לשרשרת האספקה, משום שסירב להעניק ל־Claude רישיון לכל שימוש חוקי \autocite{npr2026openai}. Maven Smart System החזירה את בית הספר כמטרה מומלצת ליום הראשון \autocite{livingstone2026maven,gizmodo2026minab}. שום דיווח שפורסם לא אומר איזה חלק היה ל־Claude בהמלצה הזאת, ואדם אחד שמכיר את העבודה של הספק אמר ל־NBC News שהמודל לא נתן ייעוץ ישיר לגבי מטרות \autocite{nbc2026targeting}.

סקירה פנימית של הפנטגון, שהממצאים המלאים שלה נשמרים בסוד כמוס, מצאה לפי הדיווחים שהכנת מטרות שארכה שעות בעימותים קודמים נדחסה לדקות, שכוח האדם שעסק בצמצום פגיעה באזרחים ירד בכתשעים אחוז, לפחות מעשרים איש, שהצוות של פיקוד המרכז של ארצות הברית ירד מעשרה לאחד, ושאף חבר בו לא בחן את האתר במינאב – תוצאה שבכירים ייחסו הן לקיצוצים והן להחלטה שלא לשתף את הצוות בתכנון ובמבצעים \autocite{Bloomberg2026KillChain}. החקירה הרשמית נפתחה במרץ, אחרי שהערכה מקדימה מצאה שארצות הברית אחראית ככל הנראה \autocite{NPR2026MinabProbe}.

במערכה באיראן ב־2026, פלטפורמת מטרות שבנה קבלן, ובתוכה מודל לשימוש כללי של ספק מסחרי, ששימשה לדירוג מטרות לפי חשיבותן האסטרטגית ולניסוח הצדקה משפטית לכל תקיפה \autocite{livingstone2026maven}, סייעה לתקוף יותר מאלף מטרות בעשרים וארבע השעות הראשונות, בערך פי עשרה ממה שהיה אפשרי לפניה \autocite{csis2026maven}, ושלושה־עשר אלף עד היום השלושים ושמונה \autocite{breakingdefense2026maven}. הפנטגון לא פרסם זמן בחינה לכל מטרה עבורה.\footnote{Maven הוא שם של תוכנית אחת שנושא יותר ממערכת אחת: תוכנית הפנטגון שעדיין נקראת Maven, ו־Maven Smart System, פלטפורמת התוכנה שצמחה ממנה \autocite{csis2026maven}. התוכנית התחילה ב־2017 כפרויקט לתיוג עצמים בווידאו שרחפנים אספו במערכה נגד המדינה האסלאמית \autocite{weisgerber2017maven,fang2019maven}. המערכת של 2026 ממזגת מערכות צבאיות ומקורות מודיעין לממשק אחד שמשמש לאיתור מטרות ולתקיפתן ולפיקוד על מבצעים \autocite{defensescoop2026transition}. הפלטפורמה היא של Palantir: ראש מערך הדיגיטל והבינה המלאכותית של הפנטגון קורא לה „מוצר התוכנה־כשירות של Palantir” \autocite{defensescoop2026transition}. העמדה הפומבית של Palantir, כפי שצוטטה מכלי שני, היא שההכרעה על חיים ומוות אינה תפקידה ושהאחריות נשארת תמיד בידי הארגון הצבאי \autocite{lee2026palantir}, והיא טוענת שהמודיעין שבבסיס היה באחריות הממשלה ושאין שום ראיה לכך ש־Maven תפקדה לא כשורה או גרמה להפצצה במינאב \autocite{Bloomberg2026KillChain}. איך עברה התוכנית ממשרד וקבלן אחד לאלה שאחריהם מופיע בנספח~\ref{sec:C.1}.}

האנשים שברשימות התקיפה האוויריות הישראליות והאמריקאיות לא יכלו להתנגד להכללתם ברשימות האלה – פסי הייצור לא נבנו כדי לקבל משוב כזה. אנדרי חוסה הרננדס רומרו נלכד בפס ייצור אחר, והוא יכול היה להתנגד. מאפר מוונצואלה, הוא התייצב במעבר הגבול בסן איסידרו באוגוסט 2024 כדי לבקש מקלט, עבר את ראיון החשש המבוסס, והיו לו עורך דין, בקשה תלויה ועומדת ומועד לדיון \autocite{LULAC2025HernandezRomero,Vega2025HernandezRomero}. עובד של קבלן שסינן אותו קרא שני קעקועים, כתרים מעל המילים „אמא” ו„אבא”, כסימני כנופיה \autocite{Palomera2025Blade}, ובמרץ 2025 הוטס לכלא באל סלוודור לפני הדיון בעניינו, בערב שבו שופט פדרלי הורה להחזיר את המטוסים \autocite{jgg2025contempt}. היה לו כל מה שאדם צריך כדי להתנגד, חוץ מהזמן. נחזור למקרה שלו בהרחבה כשנדון בפס הייצור של חטיפה והעברה בכפייה של ארצות הברית.

\textbf{מי סירב.} לתיעוד של עשרים השניות יש חצי שני, התיעוד של מי שהחזיקו מעמד מול לחץ מהסוג הזה, והוא אומר שלהרכב החדר יש חשיבות רבה. במצב הבסיס של מילגרם בערך שליש מהנבדקים הפסיקו לפני מכת החשמל המרבית. כששני עמיתים התריסו, רק אחד מכל עשרה הגיע עד הסוף, וכששני נסיינים חלקו זה על זה, הרבה פחות המשיכו מאשר במצב הבסיס \autocite{milgram1974obedience}. כרבע מהנבדקים של אש מעולם לא נכנעו לרוב הטועה \autocite{asch1956studies}. בראונינג מעריך שבין עשירית לחמישית מאנשי גדוד משטרת המילואים 101 סירבו לירי ביוזפוב, הטבח הראשון של הגדוד, או התחמקו ממנו \autocite{browning1992ordinary}. השיעור נע בין כמעט אף אחד לכמעט כולם, לפי התנאים. לאנשים שהחזיקו מעמד היו ראיות מחוץ לערוץ: עמית שהתריס, סמכות שנייה שחלקה, ובקרב המצילים שבני הזוג אוליינר חקרו, קשרים שחצו את הקו שהרדיפה סימנה \autocite{oliner1988altruistic}. או שהם החזיקו במחויבות ברמת דיוק שהלחץ לא יכול היה לגבור עליה, כמו שמחזיקים במנדטים מוסריים ובערכים מקודשים \autocite{skitka2009limits,ginges2007sacred}. או שהם התייחסו למקור כבעל אינטרס. וקבוצות התנגדו במקום שבו יחידים מבודדים כמעט אף פעם לא התנגדו \autocite{gamson1982encounters}.

אפשר להרכיב את החדר גם לטובת הצד השני, באותה מידה. לגרמניה של ויימאר היו חיי התאגדות עשירים, שאפשר לטעון שעזרו להפיל את הרפובליקה: כשמפלגות וממשלות לא נענו לטענות של אנשים, אנשים השקיעו את האנרגיה שלהם במועדונים ובאגודות שהתארגנו לפי קווים קבוצתיים, והנאצים גייסו דרך הרשתות שהאגודות האלה כבר בנו \autocite{berman1997weimar}. בעיירות שהיו בהן יותר מקהלות, מועדוני ספורט ואגודות התעמלות נרשמה הצטרפות מהירה יותר למפלגה, אבל רק במקומות שבהם השלטון היה לא יציב; בפרוסיה, שהממשלה שלה הייתה יציבה, הקשר היה קטן ולא מובהק \autocite{satyanath2017bowling}. אז מה שהיה למצילים לא היה פשוט יותר חברה. הקשרים שלהם חצו את הקו שהרדיפה סימנה. אשוטוש וארשני מצא את אותו הדבר בערי הודו: מהומות בין הינדים למוסלמים היו נדירות במקומות שבהם אגודות הביאו את שתי הקהילות לאותם ארגונים, ושכיחות במקומות שבהם לא \autocite{varshney2002ethnic}. הקשרים שמאפשרים סירוב הם אלה שתנועה פשיסטית תצטרך לנתק.

כל מה שבא מכאן והלאה עוסק בהרכבת החדר.\footnote{החדר הוא התפאורה, וכמו שהמורה שלי לתיאטרון היה אומר, התפאורה מביימת את הסצנה. מי עוד נמצא בחדר? האם אפשר לשרוד עזיבה? האם מישהו ישמע סירוב?} הרכבה היא מילה של עשייה, וחנה ארנדט הזהירה מפני יחס לפוליטיקה כאל דבר שעושים. אצל היוונים, כתבה, החוק היה פשוטו כמשמעו חומה, שבלעדיה ייתכן שיהיה יישוב, אבל לא עיר \autocite{arendt1958human}: המחוקק בנה את החומה, ומה שקרה בתוכה לא היה בידיו לעשות. האמצעים הבסיסיים שאני מצדד בהם הם חומות במובן הזה. הם מאפשרים לשרוד סירוב ודואגים שישמעו אותו; הם לא מייצרים אותו, והכוח, לפי תפיסתה, קיים רק בין אנשים הפועלים יחד ונעלם כשהם מתפזרים, וזה מה שהתיעוד של מי שסירבו מראה על סירוב. הכלים של החלקים השני והשלישי גורמים לסירוב להיחשב. המוסדות של החלק החמישי מייצרים סרבנים. מכונה שיכולה לסרב תצטרך את מה שהיה לתשעה מכל עשרה שעמית עמד לצדם: עמית בעל היסטוריה אחרת לצדה, מקום ללכת אליו, ומישהו שישמע אותה.

זה גם מה שסרבן – כולל מכונה סרבנית – בפס הייצור יכול היה להיות עבור הקצין. מכונה שסירבה הייתה על המסך כשהמקרה הגיע אליו: העמית המתריס, אם האפקט שורד את המעבר מחדר למסך, דבר שלא נבדק ושהטיית אוטומציה צופה שלא ישרוד. מכונה שהייתה מניעה בוחנים בכוח הציות שלהם לה הייתה הופכת אותם לחותמות בסימן אחר. האם הקצין עם עשרים שניות ומכונה מסרבת לצדו הוא אותו קצין כמו בתחקיר?
